\documentclass[12pt]{article}
\usepackage[utf8]{inputenc}
\usepackage[T1]{fontenc}
\usepackage{hyperref}
\title{Central Subanen Language Resources for K-12 Curriculum Development}
\author{Compiled by Claude Sonnet 5.5 (claude-sonnet-5-5) for Andr\'as Kornai}
\date{}
\begin{document}
\maketitle

\section{Introduction}
Central Subanen (ISO 639-3 syb; also Sindangan Subanun, Sinubaanen) is an Austronesian language spoken in the interior of the Zamboanga Peninsula, Mindanao, Philippines. Wikipedia notes that linguists regard the Subanen varieties as a dialect cluster. This report rests on a single web search.

\section{Primary Educational Resources}
No primer, textbook or school programme was found in the search results.

\section{Linguistic and Reference Materials}
\begin{itemize}
\item Subanen languages, Wikipedia: \url{https://en.wikipedia.org/wiki/Subanen_languages}
\item Subanon people, Wikipedia: \url{https://en.wikipedia.org/wiki/Subanon_people}
\end{itemize}
No Central Subanen dictionary or grammar was found. The tagalognow.com Subanon dictionary page appeared in results, but it was not checked for which Subanen variety it covers.

\section{Religious and Community Resources}
\begin{itemize}
\item The Bible in Central Subanen (New Testament, \emph{Kig Begu Pasad}, 1993, Wycliffe Bible Translators per the listing), Digital Bible Society: \url{https://dbs.org/bibles/SYBWBT}
\item ScriptureEarth listing for syb: \url{https://www.scriptureearth.org/00eng.php?iso=syb}
\item Joshua Project, Central Subanen in Philippines: \url{https://joshuaproject.net/people_groups/11138}
\end{itemize}

\section{Digital and Audio Resources}
Audio and text versions of the New Testament are reported available through the ScriptureEarth and Digital Bible Society listings above.

\section{Curriculum Development Context}
Search results gave no evidence of Central Subanen being taught in schools. The New Testament is the main written text found, and any curriculum would need new materials. Local Subanen community organisations and Philippine mother-tongue education programmes should be investigated further.
\end{document}